% 简要讨论初始页表条件：
%     未启用分页，但虚拟地址大多直接映射到物理地址，
%     这也是我们启用分页时的情况，
%     因为分页是在创建第一个进程后才开启的。
%   提及为何仍有 SEG_UCODE/SEG_UDATA？
%   我们是否真正解释过 %cs 的低两位的作用？
%     特别是它们与 PTE_U 的交互。
%   关于为何使用 extern char[] 的边注
\chapter{页表}
\label{CH:MEM}

页表是操作系统为每个进程提供其各自私有地址空间和内存的最常用机制。页表决定了内存地址的含义，以及物理内存的哪些部分可以被访问。它们使得 xv6 能够隔离不同进程的地址空间，并将它们多路复用到一个单一的物理内存上。页表提供了一层间接性，允许操作系统执行许多有用的技巧。Xv6 实现了其中几个技巧：将同一块内存（跳板页）映射到多个地址空间中，用未映射的页来保护内核和用户栈，以及惰性分配用户堆内存。本章其余部分将解释 RISC-V 硬件提供的页表以及 xv6 如何使用它们。

\section{分页硬件}
提醒一下，RISC-V 指令（包括用户和内核指令）操作的是虚拟地址。机器的 RAM，或者说物理内存，是通过物理地址来索引的。RISC-V 页表硬件通过将每个虚拟地址映射到一个物理地址，将这两种地址连接起来。

\begin{figure}[t]
\center
\includegraphics[scale=0.5]{fig/riscv_address.pdf}
\caption{将虚拟地址映射到物理地址的扁平页表的抽象视图。}
\label{fig:riscv_address}
\end{figure}

Xv6 使用 RISC-V 的 Sv39 模式，这意味着 64 位虚拟地址中只使用低 39 位；高 25 位不被使用。在这种 Sv39 配置下，一个 RISC-V 页表逻辑上是一个包含 $2^{27}$（134,217,728）个页表项（PTE）的数组。每个 PTE 包含一个 44 位的物理页号（PPN）和一些标志位。分页硬件通过使用 39 位中的高 27 位作为页表索引来查找 PTE，然后构造一个 56 位的物理地址，其高 44 位来自 PTE 中的 PPN，低 12 位直接复制自原始虚拟地址。图~\ref{fig:riscv_address} 以页表作为 PTE 的简单数组的逻辑视图展示了这个过程（RISC-V 页表实际上是一个树；更完整的情况见图~\ref{fig:riscv_pagetable}）。页表以 4096（$2^{12}$）字节的对齐块为粒度，让操作系统能够控制虚拟地址到物理地址的转换。这样的一个块被称为一页。

\begin{figure}[t]
\center
\includegraphics[scale=0.5]{fig/riscv_pagetable.pdf}
\caption{RISC-V 地址转换细节。}
\label{fig:riscv_pagetable}
\end{figure}

RISC-V 的设计为虚拟地址和物理地址的扩展预留了空间。如果需要更大的虚拟地址空间，RISC-V 支持 Sv48 模式，使用 48 位虚拟地址~\cite{riscv:priv}。物理地址也有增长空间：PTE 格式中为物理页号预留了再增加 10 位的空间。RISC-V 的设计者根据技术预测选择了地址大小。$2^{48}$ 字节是 262,144 GB，这比目前任何应用可能使用的用户虚拟地址空间都要大得多。$2^{56}$ 字节的物理地址空间是 65,536 TB，远远超过当前任何计算机所能配备的内存容量。

如图~\ref{fig:riscv_pagetable} 所示，RISC-V CPU 页表作为三层树存储在物理内存中。树的根是一个 4096 字节的页表页，包含 512 个 PTE，这些 PTE 包含树中下一级页表页的物理地址。这些页表中的每一个又包含 512 个 PTE，用于树的最后一级。分页硬件使用 27 位地址中的高 9 位在根页表页中选择一个 PTE，中间 9 位在树的下一个级别的页表页中选择一个 PTE，低 9 位则选择最终的 PTE。（在 Sv48 RISC-V 中，页表有四层，虚拟地址的第 39 至 47 位用于索引顶层。）

如果转换一个地址所需的三个 PTE 中有任何一个不存在，分页硬件就会引发一个页错误异常，由内核来处理该页错误（参见第~\ref{CH:TRAP} 章和第~\ref{CH:PGFAULTS} 章）。

与图~\ref{fig:riscv_address} 的单级设计相比，图~\ref{fig:riscv_pagetable} 的三层结构提供了一种内存高效的方式来记录 PTE。在常见的场景中，大范围的虚拟地址没有映射，三层结构可以省略整个页目录。例如，如果一个应用只使用从地址零开始的少数几页，那么顶层页目录的第 1 至 511 项是无效的，内核就不必为这 511 个中间页目录分配页。此外，内核也不必为这 511 个中间页目录分配底层页目录的页。因此，在这个例子中，三层设计节省了 511 页用于中间页目录和 $511\times512$ 页用于底层页目录。

尽管 CPU 在执行加载或存储指令时会在硬件中遍历三层结构，但三层的潜在缺点是 CPU 必须从内存中加载三个 PTE 才能将加载/存储指令中的虚拟地址转换为物理地址。为了避免从物理内存加载 PTE 的开销，RISC-V CPU 会将页表项缓存到转址旁路缓存（TLB）中。

每个 PTE 包含标志位，用于告诉分页硬件如何使用相关联的虚拟地址。
\indexcode{PTE_V}
指示 PTE 是否存在：如果未设置，则访问该页会引发页错误（即不允许访问）。
\indexcode{PTE_R}
控制是否允许指令读取该页。
\indexcode{PTE_W}
控制是否允许指令写入该页。
\indexcode{PTE_X}
控制 CPU 是否可以将该页的内容解释为指令并执行。
\indexcode{PTE_U}
控制用户模式下的指令是否允许访问该页；
如果 \indexcode{PTE_U} 未设置，则 PTE 只能在监督者模式下使用。
图~\ref{fig:riscv_pagetable} 显示了这些标志位在 PTE 中的位置。
这些标志位以及所有其他与页硬件相关的结构定义在
\fileref{kernel/riscv.h} 中。

要告诉 CPU 使用某个页表，内核必须将根页表页的物理地址写入
\texttt{satp}\index{satp@\lstinline{satp}} 寄存器。
CPU 将使用其 \texttt{satp} 所指向的页表来转换随后所有指令生成的地址。
每个 CPU 都有自己的 \texttt{satp}，这样不同的 CPU 就可以运行不同的进程，每个进程都有自己私有的地址空间，由各自的页表描述。

从内核的角度来看，页表是存储在内存中的数据，内核使用与处理任何树形数据结构类似的代码来创建和修改页表。

关于本书中使用的术语有几点说明。\textit{物理内存} 指的是 RAM 中的存储单元。一个字节的物理内存有一个地址，称为 \textit{物理地址}。解引用地址的指令（如加载、存储、跳转和函数调用）只使用虚拟地址，分页硬件将这些虚拟地址转换为物理地址，然后发送到 RAM 硬件进行读写存储。\textit{地址空间} 是在给定页表中有效的虚拟地址集合；每个 xv6 进程都有一个独立的用户地址空间，xv6 内核也有自己的地址空间。\textit{用户内存} 指的是进程的用户地址空间加上页表允许该进程访问的物理内存。\textit{虚拟内存} 指的是与管理页表以及使用它们实现隔离等目标相关的思想和技术。

\begin{figure}[h]
\centering
 \includegraphics[scale=0.65]{fig/xv6_layout.pdf}
\caption{左图：xv6 的内核虚拟地址空间。
{\sf \small{RWX}}
指的是 PTE 的读、写和执行权限。
右图：xv6 期望看到的 RISC-V 物理地址空间。}
\label{fig:xv6_layout}
\end{figure}

\section{内核地址空间}
启动时，xv6 会创建一个描述内核地址空间的单一页表。内核配置其地址空间的布局，以便在可预测的虚拟地址上访问物理内存和各种硬件资源。图~\ref{fig:xv6_layout} 展示了这种布局如何将内核虚拟地址映射到物理地址。文件 \fileref{kernel/memlayout.h} 声明了 xv6 内核内存布局的常量。

QEMU 模拟的计算机包括从物理地址 \texttt{0x80000000} 开始，至少持续到 \texttt{0x88000000} 的 RAM（物理内存），xv6 将其称为 \texttt{PHYSTOP}。QEMU 模拟还包括 I/O 设备，例如磁盘接口。QEMU 将设备接口作为内存映射的控制寄存器暴露给软件，这些寄存器位于物理地址空间的 \texttt{0x80000000} 以下。内核可以通过读写这些特殊的物理地址来与设备交互；这样的读写操作是与设备硬件通信，而不是与 RAM 通信。第~\ref{CH:TRAP} 章解释了 xv6 如何与设备交互。

内核将所有物理 RAM 和设备寄存器映射到与物理地址相等的虚拟地址上。这被称为“直接映射”，它允许内核只需加载或存储虚拟地址 $x$ 就可以读写物理地址 $x$。内核代码本身在虚拟地址空间和物理内存中都位于 \lstinline{KERNBASE=0x80000000}。当 \lstinline{kfork} \lineref{kernel/proc.c:/^kfork/} 为子进程分配用户内存时，分配器返回该内存的物理地址；\lstinline{kfork} 在将父进程的用户内存复制到子进程时，直接将该地址用作虚拟地址。

有几个内核虚拟地址不是直接映射的：

\begin{itemize}
  
\item 跳板页。它被映射在虚拟地址空间的顶部；用户页表也有这个相同的映射。第~\ref{CH:TRAP} 章讨论了跳板页的作用，但在这里我们看到了页表的一个有趣用例：一个物理页（存放跳板代码）在内核的虚拟地址空间中被映射了两次：一次在虚拟地址空间的顶部，一次通过直接映射。

\item 内核栈页。每个进程都有自己的内核栈，它被映射到一个高地址的内核虚拟地址，这样 xv6 就可以在其下方留出一个未映射的防护页。防护页的 PTE 是无效的（即 \lstinline{PTE_V} 未设置），这样如果内核栈溢出，很可能会导致页错误，然后内核 panic。如果没有防护页，溢出的栈会覆盖其他内核内存，导致错误操作。相比之下，panic 崩溃更可取。

\end{itemize}

虽然内核通过高内存映射使用其栈，但每个栈也可以通过直接映射地址被内核访问。另一种设计可能只使用直接映射，并在直接映射地址上使用栈。然而，在这种安排下，提供防护页需要取消映射那些本来指向物理内存的虚拟地址，这会使这些地址难以使用。

内核将跳板页和内核文本的页映射为具有权限 \lstinline{PTE_R} 和 \lstinline{PTE_X}，但没有 \lstinline{PTE_W}。内核将其它页映射为具有权限 \lstinline{PTE_R} 和 \lstinline{PTE_W}，但没有 \lstinline{PTE_X}。防护页的映射是无效的。这些受限权限的目的是帮助捕获以意外方式访问页的内核错误，例如内核代码意外地试图写入内核指令。

内核创建一个单一的内核页表，所有 CPU 在内核中执行时都使用它。xv6 在初始创建后不会修改内核页表。

\section{代码：创建地址空间}
请在继续阅读之前，通读 {\tt kernel/vm.c} 到 {\tt mappages()} 结束的部分。

xv6 中用于操作地址空间和页表的大部分代码位于 {\tt vm.c} \fileref{kernel/vm.c} 中。核心数据结构是 {\tt pagetable\_t}，它实际上是一个指向 RISC-V 根页表页的指针；{\tt pagetable\_t} 可以是内核页表，也可以是每个进程的页表之一。核心函数是 {\tt walk}（用于查找虚拟地址对应的 PTE）和 {\tt mappages}（用于为新映射安装 PTE）。以 {\tt kvm} 开头的函数操作内核页表；以 {\tt uvm} 开头的函数操作用户页表；其他函数则两者都用。{\tt copyout} 和 {\tt copyin} 用于在系统调用参数提供的用户虚拟地址之间复制数据；它们位于 {\tt vm.c} 中，因为它们需要显式转换这些地址以找到对应的物理内存。

在启动序列的早期，\indexcode{main} 调用 \indexcode{kvminit} \lineref{kernel/vm.c:/^kvminit/}，通过 \indexcode{kvmmake} \lineref{kernel/vm.c:/^kvmmake/} 创建内核页表。这个调用发生在 xv6 为 RISC-V 启用分页之前，因此地址直接指向物理内存。\lstinline{kvmmake} 首先分配一个物理内存页来存放根页表页。然后它调用 \indexcode{kvmmap} 来安装内核所需的转换。这些转换包括内核的指令和数据、直到 \indexcode{PHYSTOP} 的物理内存，以及实际上是设备的内存范围。\indexcode{proc_mapstacks} \lineref{kernel/proc.c:/^proc_mapstacks/} 为每个进程分配一个内核栈。它调用 \lstinline{kvmmap} 将每个栈映射到由 \lstinline{KSTACK} 生成的虚拟地址上，该映射为无效的栈防护页留出了空间。

\indexcode{kvmmap} \lineref{kernel/vm.c:/^kvmmap/} 调用 \indexcode{mappages} \lineref{kernel/vm.c:/^mappages/}，后者将一段虚拟地址范围映射到对应的物理地址范围，并将映射安装到页表中。它按页间隔对范围内的每个虚拟地址分别执行此操作。对于每个要映射的虚拟地址，\lstinline{mappages} 调用 \indexcode{walk} 来查找该地址对应的 PTE 的地址。然后它初始化 PTE，使其包含相关的物理页号、所需的权限（\lstinline{PTE_W}、\lstinline{PTE_X} 和/或 \lstinline{PTE_R}）以及 \lstinline{PTE_V} 来标记 PTE 为有效 \lineref{kernel/vm.c:/perm...PTE_V/}。

\indexcode{walk} \lineref{kernel/vm.c:/^walk/} 模仿 RISC-V 分页硬件查找虚拟地址对应 PTE 的过程（见图~\ref{fig:riscv_pagetable}）。\lstinline{walk} 一次下降一层页表树，使用虚拟地址每层的 9 位来索引相关的页目录页。在每一层，它要么找到指向下一层页目录页的 PTE，要么找到最终页的 PTE \lineref{kernel/vm.c:/pte.=..pagetable/}。如果第一层或第二层页目录页中的某个 PTE 无效，那么所需的目录页尚未分配；如果设置了 \lstinline{alloc} 参数，\lstinline{walk} 会分配一个新的页表页并将其物理地址放入该 PTE。它返回树中最底层 PTE 的地址 \lineref{kernel/vm.c:/return..pagetable/}。

上述代码依赖于物理内存被直接映射到内核虚拟地址空间。例如，当 \lstinline{walk} 下降页表的层级时，它会从 PTE 中取出下一级页表页的（物理）地址 \lineref{kernel/vm.c:/pagetable.=..pa.*E2P/}，然后将该地址作为虚拟地址来获取下一级的 PTE \lineref{kernel/vm.c:/t..pte.=..paget/}。

在每个 CPU 上，\indexcode{main} 调用 \indexcode{kvminithart} \lineref{kernel/vm.c:/^kvminithart/} 来安装内核页表，即将根页表页的物理地址放入 CPU 的 \texttt{satp} 寄存器。此后，CPU 使用内核页表转换地址。内核继续正确执行，因为内核页表是直接映射的，所以在此更改前后，地址都指向 RAM 中的相同位置。

每个 RISC-V CPU 都会将页表项缓存到转址旁路缓存（TLB）中，当 xv6 更改页表时，它必须告诉 CPU 使相应的缓存 TLB 项失效。如果不这样做，那么稍后 TLB 可能会使用旧的缓存映射，指向一个在此期间已分配给另一个进程的物理页，结果可能导致一个进程能够篡改另一个进程的内存。RISC-V 有一条指令 \indexcode{sfence.vma} 用于刷新当前 CPU 的 TLB。Xv6 在重新加载 \texttt{satp} 寄存器之后，在 {\tt kvminithart} 中执行 {\tt sfence.vma}，并在 \lstinline{uservec} 和 \lstinline{userret} 的跳板代码中也执行该指令。

在更改 \texttt{satp} 之前也有必要执行 \texttt{sfence.vma}，以等待所有未完成的加载和存储完成。这种等待确保了先前对页表的更新已经完成，并确保先前的加载和存储使用的是旧的页表，而不是新的。

\section{物理内存分配}
内核必须在运行时为页表、用户内存、内核栈和管道缓冲区分配和释放物理内存。

Xv6 使用内核结束之后到 \indexcode{PHYSTOP} 之间的物理内存进行运行时分配。它一次分配和释放整个 4096 字节的页。它通过将这些页本身穿成一个链表来跟踪哪些页是空闲的。分配操作包括从链表中移除一页；释放操作包括将释放的页添加到链表中。

请阅读 {\tt kernel/kalloc.c}。

\section{代码：物理内存分配器}
分配器位于 {\tt kalloc.c} \fileref{kernel/kalloc.c} 中。分配器的数据结构是一个由可供分配的物理内存页构成的空闲链表。每个空闲页的“next”指针存放在一个 \indexcode{struct run} \lineref{kernel/kalloc.c:/^struct.run/} 中。分配器将每个空闲页的 \lstinline{run} 结构存储在该空闲页本身中，因为当页空闲时，里面没有存储其他东西。空闲链表由一个自旋锁保护 \linerefs{kernel/kalloc.c:/^struct.{/,/}/}。链表和锁被封装在一个结构体中，以明确该锁保护的是结构体中的字段。现在，请忽略锁以及对 \lstinline{acquire} 和 \lstinline{release} 的调用；第~\ref{CH:LOCK} 章将详细讨论锁。

函数 \indexcode{main} 调用 \indexcode{kinit} 来初始化分配器 \lineref{kernel/kalloc.c:/^kinit/}。\lstinline{kinit} 初始化空闲链表，使其包含内核结束之后到 {\tt PHYSTOP} 之间的每一页物理 RAM。Xv6 本应通过解析硬件提供的配置信息来确定实际有多少物理内存。但 xv6 假定机器有 128 兆字节的 RAM。\lstinline{kinit} 调用 \indexcode{freerange}，通过逐页调用 \indexcode{kfree} 将内存添加到空闲链表。PTE 只能引用在 4096 字节边界上对齐（即为 4096 的倍数）的物理地址，因此 \lstinline{freerange} 使用 \indexcode{PGROUNDUP} 来确保它只释放对齐的物理地址。分配器开始时没有内存；这些对 \lstinline{kfree} 的调用给它提供了一些内存进行管理。

分配器有时将地址视为整数以对其执行算术运算（例如，在 \lstinline{freerange} 中遍历所有页），有时又将地址用作指针来读写内存（例如，操作存储在每个页中的 \lstinline{run} 结构）；地址的这种双重用途是分配器代码中包含大量 C 类型转换的主要原因。

函数 \lstinline{kfree} \lineref{kernel/kalloc.c:/^kfree/} 首先将正在释放的内存中的每个字节设置为 1。这将导致在释放后使用内存的代码（“悬空引用”）读到的是垃圾值而不是旧的有效内容；希望这能使此类代码更快地出错。然后 \lstinline{kfree} 将该页添加到空闲链表头部：它将 \lstinline{pa} 强制转换为指向 \lstinline{struct run} 的指针，将旧的空闲链表头部记录在 \lstinline{r->next} 中，并将空闲链表头部设置为 \lstinline{r}。\indexcode{kalloc} 移除并返回空闲链表中的第一个元素。

\section{进程地址空间}
每个进程都有自己的页表，当 xv6 在进程间切换时，它也会更改页表。图~\ref{fig:processlayout} 比图~\ref{fig:as} 更详细地展示了一个进程的地址空间。进程的用户地址空间从零开始，原则上结束于 \texttt{MAXVA}（{\tt 0x4000000000}）\lineref{kernel/riscv.h:/MAXVA/}，但实际上只有一小部分被映射到物理内存。

一个进程的地址空间由若干页组成，这些页包含程序的文本（xv6 以 \lstinline{PTE_R}、\lstinline{PTE_X} 和 \lstinline{PTE_U} 权限映射）、包含程序预初始化数据的页、一页用于栈，以及用于堆的页。Xv6 以 \lstinline{PTE_R}、\lstinline{PTE_W} 和 \lstinline{PTE_U} 权限映射数据、栈和堆。

在用户地址空间内使用权限是加固用户进程的常用技术。如果文本以 \lstinline{PTE_W} 映射，那么进程可能会意外修改自己的程序；例如，编程错误可能导致程序向空指针写入，修改地址 0 处的指令，然后继续运行，可能会造成更大的破坏。为了立即检测到此类错误，xv6 在映射文本时不带 \lstinline{PTE_W}；如果程序意外尝试写入地址 0，硬件将拒绝执行该写入并引发页错误（参见第~\ref{CH:TRAP} 章）。然后内核会终止该进程并打印出一条信息性消息，以便开发人员追踪问题。

类似地，通过在没有 \lstinline{PTE_X} 的情况下映射数据，用户程序无法意外跳转到程序数据中的地址并开始在那里执行。

在现实世界中，通过仔细设置权限来加固进程也有助于防御安全攻击。攻击者可能向程序（例如，Web 服务器）提供精心构造的输入，以触发程序中的错误，希望将该错误转化为漏洞~\cite{aleph:smashing}。仔细设置权限以及其他技术（例如随机化用户地址空间的布局）会使此类攻击更加困难。

栈是一个单独的页，图中显示了由 {\tt exec} 系统调用创建的初始内容。包含命令行参数的字符串以及指向它们的指针数组位于栈的最顶部。就在其下方是一些值，这些值使得程序可以像刚刚调用了函数 \lstinline{main(argc}, \lstinline{argv)} 一样从 \lstinline{main} 开始运行。

为了检测用户栈溢出分配的栈内存，xv6 通过在栈正下方清除 \lstinline{PTE_U} 标志位来放置一个不可访问的防护页。如果用户栈溢出并且进程试图使用栈下方的地址，硬件将产生页错误异常，因为防护页在用户模式下运行的程序是无法访问的。在实际的操作系统中，当用户栈溢出时，可能会自动为其分配更多内存。

这里我们看到了页表使用的几个很好的例子。第一，不同进程的页表将用户地址转换为不同物理内存页，因此每个进程都拥有私有的用户内存。第二，每个进程看到自己的内存具有从零开始的连续虚拟地址，而进程的物理内存可以是非连续的。第三，内核在用户地址空间的顶部映射了一个包含跳板代码的页（不带 \lstinline{PTE_U}），因此一个物理内存页出现在所有地址空间中，但只能被内核使用。

\begin{figure}[t]
\center
\includegraphics[scale=0.5]{fig/processlayout.png}
\caption{一个进程的用户地址空间，及其初始栈。}
\label{fig:processlayout}
\end{figure}

\section{代码：exec}
请在继续之前阅读 {\tt kernel/exec.c} 以及从 {\tt uvmcreate()} 开始的 {\tt kernel/vm.c}。

\lstinline{exec} 是一个系统调用，它用从文件（称为二进制文件或可执行文件）中读取的数据替换进程的用户地址空间。二进制文件通常是编译器和链接器的输出，包含机器指令和程序数据。\lstinline{kexec} \lineref{kernel/exec.c:/^kexec/} 是 {\tt exec} 的内核内部实现，它使用 \indexcode{namei} \lineref{kernel/exec.c:/namei/} 打开指定路径的二进制文件，这将在第~\ref{CH:FS} 章中解释。然后，它读取 ELF 头。Xv6 二进制文件采用广泛使用的 ELF 格式，定义在 \fileref{kernel/elf.h} 中。一个 ELF 二进制文件由一个 ELF 头 \indexcode{struct elfhdr} \lineref{kernel/elf.h:/^struct.elfhdr/} 以及随后的一系列程序段头 \lstinline{struct proghdr} \lineref{kernel/elf.h:/^struct.proghdr/} 组成。每个 \lstinline{proghdr} 描述了应用程序中必须加载到内存的一个段；xv6 程序有两个程序段头：一个用于指令，一个用于数据。

第一步是快速检查该文件可能包含一个 ELF 二进制文件。一个 ELF 二进制文件以四个字节的“幻数” \lstinline{0x7F}、\lstinline{`E`}、\lstinline{`L`}、\lstinline{`F`} 或 \indexcode{ELF_MAGIC} \lineref{kernel/elf.h:/ELF_MAGIC/} 开头。如果 ELF 头具有正确的幻数，\lstinline{kexec} 就假定该二进制文件格式良好。

\lstinline{kexec} 使用 \indexcode{proc_pagetable} \lineref{kernel/exec.c:/proc_pagetable/} 分配一个没有用户映射的新页表，使用 \indexcode{uvmalloc} \lineref{kernel/exec.c:/uvmalloc/} 为每个 ELF 段分配内存，并使用 \indexcode{loadseg} \lineref{kernel/exec.c:/loadseg/} 将每个段加载到内存中。\lstinline{loadseg} 使用 \indexcode{walkaddr} 找到分配的内存的物理地址，以便写入 ELF 段的每一页，并使用 \indexcode{readi} 从文件中读取。

第一个用户程序 \indexcode{/init} 的程序段头如下所示：
\begin{footnotesize}
\begin{verbatim}
# objdump -p user/_init

user/_init:     file format elf64-little

Program Header:
0x70000003 off    0x0000000000006bb0 vaddr 0x0000000000000000
                                       paddr 0x0000000000000000 align 2**0
         filesz 0x000000000000004a memsz 0x0000000000000000 flags r--
    LOAD off    0x0000000000001000 vaddr 0x0000000000000000
                                       paddr 0x0000000000000000 align 2**12
         filesz 0x0000000000001000 memsz 0x0000000000001000 flags r-x
    LOAD off    0x0000000000002000 vaddr 0x0000000000001000
                                       paddr 0x0000000000001000 align 2**12
         filesz 0x0000000000000010 memsz 0x0000000000000030 flags rw-
   STACK off    0x0000000000000000 vaddr 0x0000000000000000
                                       paddr 0x0000000000000000 align 2**4
         filesz 0x0000000000000000 memsz 0x0000000000000000 flags rw-
\end{verbatim}
\end{footnotesize}

我们看到文本段应该从文件偏移量 0x1000 处加载到内存中的虚拟地址 0x0（不带写权限）。我们还看到数据段应该加载到地址 0x1000（在页边界上），且不带可执行权限。

程序段头的 \lstinline{filesz} 可能小于 \lstinline{memsz}，表示它们之间的间隙应该用零填充（用于 C 全局变量），而不是从文件中读取。对于 \lstinline{/init}，数据段的 \lstinline{filesz} 是 0x10 字节，\lstinline{memsz} 是 0x30 字节，因此 \indexcode{uvmalloc} 分配足够的物理内存来容纳 0x30 字节，但仅从文件 \lstinline{/init} 中读取 0x10 字节。

现在，\indexcode{kexec} 分配并初始化用户栈。它只分配一个栈页。\lstinline{kexec} 将参数字符串逐个复制到栈顶，并将指向它们的指针记录在 \indexcode{ustack} 中。它在最终将传递给 \lstinline{main} 的 \indexcode{argv} 列表末尾放置一个空指针。传递给 \lstinline{main} 的 \indexcode{argc} 和 \lstinline{argv} 值通过系统调用返回路径传递：\lstinline{argc} 通过系统调用返回值传递，该值存放在 {\tt a0} 中，而 \lstinline{argv} 通过进程陷阱帧的 {\tt a1} 条目传递。

\lstinline{kexec} 在栈页的正下方放置一个不可访问的页，这样尝试使用超过一页的程序就会发生故障。这个不可访问的页还允许 \lstinline{kexec} 处理过大的参数；在这种情况下，\lstinline{kexec} 用来将参数复制到栈的 \indexcode{copyout} \lineref{kernel/vm.c:/^copyout/} 函数会发现目标页不可访问，并返回 -1。

在准备新内存映像的过程中，如果 \lstinline{kexec} 检测到错误（例如无效的程序段），它会跳转到标签 \lstinline{bad}，释放新映像，并返回 -1。\lstinline{kexec} 必须等到确定系统调用会成功之后才能释放旧映像：如果旧映像被释放了，系统调用就无法向其返回 -1。\lstinline{kexec} 中唯一的错误情况发生在创建新映像的过程中。一旦映像构建完成，\lstinline{kexec} 就可以切换到新页表 \lineref{kernel/exec.c:/pagetable.=.pagetable/} 并释放旧页表 \lineref{kernel/exec.c:/proc_freepagetable/}。

\lstinline{exec} 系统调用将字节从 ELF 文件加载到 ELF 文件指定的地址处的内存中。用户或进程可以将任何他们想要的地址放入 ELF 文件中。因此 \lstinline{exec} 是有风险的，因为 ELF 文件中的地址可能意外或故意地指向内核。对于一个不谨慎的内核，后果可能从崩溃到恶意破坏内核隔离机制（即安全漏洞）不等。Xv6 执行了许多检查来避免这些风险。例如，\lstinline{if(ph.vaddr + ph.memsz < ph.vaddr)} 检查总和是否会溢出 64 位整数。危险在于，用户可以构造一个 ELF 二进制文件，其中 \lstinline{ph.vaddr} 指向一个用户选择的地址，而 \lstinline{ph.memsz} 足够大，使得总和溢出到 0x1000，这看起来像一个有效的值。在旧版本的 xv6 中，用户地址空间也包含内核（但在用户模式下不可读/写），用户可以选择一个对应内核内存的地址，从而将数据从 ELF 二进制文件复制到内核中。在 RISC-V 版本的 xv6 中，这种情况不会发生，因为内核有自己的独立页表；\lstinline{loadseg} 加载到进程的页表中，而不是内核的页表。

内核开发者很容易遗漏关键的检查，而现实世界的内核有着悠久的缺失检查的历史，用户程序可以利用这些缺失来获得内核特权。很可能 xv6 并没有完全做好验证提供给内核的用户级数据的工作，恶意用户程序可能能够利用这一点来绕过 xv6 的隔离机制。

\section{现实世界}
与大多数操作系统一样，xv6 使用分页硬件进行内存保护和映射。大多数操作系统通过结合分页和页错误异常，对分页的使用远比 xv6 复杂，我们将在第~\ref{CH:TRAP} 章讨论这一点。

Xv6 之所以简化，是因为内核使用了虚拟地址和物理地址之间的直接映射，并且假设物理 RAM 位于地址 0x80000000（内核期望被加载的位置）。这在 QEMU 上可以工作，但在真实硬件上，这是一个坏主意；真实硬件会将 RAM 和设备放置在不可预测的物理地址上，因此（例如）在 xv6 期望能够存储内核的地址 0x80000000 处可能根本没有 RAM。更严谨的内核设计会利用页表将任意的硬件物理内存布局转换为可预测的内核虚拟地址布局。

RISC-V 支持物理地址级别的保护，但 xv6 没有使用该特性。

在具有大量内存的机器上，使用 RISC-V 对“大页”的支持可能是有意义的。当物理内存较小时，小页是有意义的，因为它允许以细粒度进行分配和交换到磁盘。例如，如果一个程序只使用 8 KB 内存，给它一个完整的 4 MB 大页物理内存是浪费的。大页在具有大量 RAM 的机器上更有意义，并且可以减少页表操作的开销。

为了避免在更改页表时刷新整个 TLB，RISC-V CPU 可能支持地址空间标识符（ASID）~\cite{riscv:priv}。这样内核就可以只刷新特定地址空间的 TLB 项。Xv6 没有使用这个特性。

xv6 内核缺乏一个能够提供小对象内存的类似于 {\tt malloc} 的分配器，这妨碍了内核使用需要动态分配的复杂数据结构。一个更复杂的内核可能会分配许多不同大小的小块，而不仅仅是像 xv6 那样只分配 4096 字节的块；一个真正的内核分配器需要同时处理小分配和大分配。

内存分配是一个永恒的热门话题，基本问题在于有效利用有限的内存和为未知的未来请求做准备~\cite{knuth}。如今人们更关心速度而不是空间效率。

\section{练习}
\begin{enumerate}
  
\item 解析 RISC-V 的设备树以找出计算机拥有的物理内存数量。

\item 函数 {\tt copyin} 和 {\tt copyinstr} 在软件中遍历用户页表。设置内核页表，使得内核映射了用户程序，并且 {\tt copyin} 和 {\tt copyinstr} 可以使用 {\tt memcpy} 将系统调用参数复制到内核空间，依靠硬件来执行页表遍历。

\item 修改 xv6，为内核使用大页。

\item Unix 的 \lstinline{exec} 实现传统上包含对 shell 脚本的特殊处理。如果要执行的文件以文本 \lstinline{#!} 开头，则第一行被视为用于解释该文件的程序。例如，如果调用 \lstinline{exec} 来运行 \lstinline{myprog} \lstinline{arg1}，并且 \lstinline{myprog} 的第一行是 \lstinline{#!/interp}，那么 \lstinline{exec} 会运行 \lstinline{/interp}，命令行参数为 \lstinline{/interp} \lstinline{myprog} \lstinline{arg1}。在 xv6 中实现对这一约定的支持。

\item 为内核实现地址空间布局随机化。

\end{enumerate}
